\documentclass{article}
\usepackage[utf8]{inputenc}
\usepackage[T1]{fontenc}
\usepackage{geometry}
\usepackage{graphicx}
\usepackage{float}
\usepackage{hyperref}
\usepackage[english,indonesian]{babel}

\title{Optimasi Pemrosesan Template LaTeX Menggunakan Pendekatan Berbasis AI}
\author{Academic Support Agent}
\date{\today}

\begin{document}
\maketitle

\begin{abstract}

\end{abstract}

\chapter{PENDAHULUAN}


\chapter{TINJAUAN PUSTAKA}


% Original content (commented out):
% \documentclass[12pt]{article}
\usepackage{thianswersheetAA}

\usepackage{ifxetex}
\ifxetex
\usepackage{xltxtra} 
\setmainfont[
Path=Charis Font/,
BoldFont       = CharisSILB.ttf,
ItalicFont     = CharisSILI.ttf,
BoldItalicFont = CharisSILBI.ttf
]{CharisSILR.ttf}
\else
\usepackage{charter}
\usepackage[utf8]{vntex}
\fi
\usepackage[
marginparsep=6pt,
top=1.5cm, 
bottom=2cm,
left=1cm,
right=1cm,
bindingoffset=0.5cm,
footskip=0.75cm,
twoside,
a4paper]{geometry}
\usepackage{fancyhdr}
\pagestyle{fancy}
\fancyhf{}
\renewcommand{\headrulewidth}{0pt}
\tentruonghoc{ĐH DUY TÂN}
\tendonvi{Khoa: MT \& KHTN}
\tenkithi{KÌ THI KẾT THÚC HỌC PHẦN}
\tendethiso{1}
\dauphantachgiuaSodevaMade{.}
\soluongphienbanmadethi{4} %4 là số lượng mã đề thi khác nhau được in ra
\chuphienbanmadethi{\Alph}%{\alph}%{\Roman}

% THIẾT LẬP TỔNG ĐIỂM CỦA CÁC PHẦN
\tongdiemcauhoiTracNghiem{4.0}
\tongdiemcauhoiNgan{3.0}
\tongdiemcauhoiTuLuan{3.0}

% THIẾT LẬP CÁC CÂU HỎI TNKQ
\socauhoiTracNghiem{15}
\socauhoiTracNghiemtrenCot{20}% chia số lượng câu hỏi Trắc Nghiệm trong 1 cột
\sophuongantraloiTracNghiem{4}

% THIẾT LẬP CÁC CÂU HỎI NGẮN
\socauhoiNgan{6}  % số lượng câu hỏi ngắn
\sodongmoicauhoiNgan{3}


\mauchucauhoi{red} % chọn màu cho chữ câu hỏi
\tenhocphan{Lí thuyết xác suất \& thống kê toán}
\tenmamon{STA151}
\thoigianlambai{75}

\begin{document}

\ingiaylambaiCauHoiCotAA{false}{blue}% số câu hỏi được xếp theo Cột, bạn có thể thay true bởi false để xem hiển thị

\phantieudeTrangAnswerAA

\answerCauHoiCotAA{
1;2;4;3;2;
3;1;2;0;0;
1;2;4;3;2
}{blue}{1}
\hfill
\answerCauHoiCotAA{
1;2;4;3;2;
3;1;2;2;4;
1;2;4;3;2
}{blue}{2}%
\hfill
\answerCauHoiCotAA{
3;2;4;1;2;
0;0;1;2;4;
4;2;3;1;4
}{blue}{3}%
\hfill
\answerCauHoiCotAA{
2;4;3;2;3;
1;2;2;1;3;
2;2;1;2;4
}{blue}{4}%

\end{document}


\end{document}